\section{The free stress-tensor trace}
\label{app:gten}

This appendix gives the calculation behind theorem~\ref{thm:gten} in full, and
records how it is verified.

\paragraph{The stress tensor.}
For a $p$-form with field strength $F$ of rank $p$ in $d$ dimensions, the free
stress tensor takes the form
\begin{equation}
T_{\mu\nu}
 = \frac{1}{(p-1)!}\,F_{\mu\alpha_2\dots\alpha_p}
   F_\nu{}^{\alpha_2\dots\alpha_p}
 - c\,\eta_{\mu\nu}\,\langle F,F\rangle,
\label{eq:stress}
\end{equation}
where $c$ is an improvement coefficient. Different conventions in the literature
fix $c$ differently; the point of what follows is that the conclusion does not
depend on it.

\paragraph{The trace.}
Taking the trace of \eqref{eq:stress} with $\eta^{\mu\nu}$, the first term gives
$\langle F,F\rangle$ up to normalisation and the second gives
$-cd\,\langle F,F\rangle$, so
\begin{equation}
T^\mu{}_\mu = (1 - cd)\,\langle F,F\rangle .
\label{eq:trace}
\end{equation}
Whatever $c$ is, the trace is a multiple of the single scalar
$\langle F,F\rangle$. There is no other quadratic Lorentz scalar available: the
only way to contract two copies of a five-form into a scalar is the full
contraction.

\paragraph{The vanishing.}
Here $p = 5$ and $d = 10$, so $F$ is a middle-degree form of \emph{odd} degree in
an even-dimensional space. Consider the top form $F \wedge F$. Interchanging the
two factors of a $q$-form costs a sign $(-1)^{q^2} = (-1)^q$, which for $q = 5$
is $-1$; hence $F\wedge F = -F\wedge F$ and $F \wedge F \equiv 0$ identically,
for every five-form whatsoever.

Now $\langle F, \star F\rangle\,\mathrm{vol} \propto F \wedge F$, so
$\langle F, \star F\rangle = 0$. On either eigenspace of the star,
$\star F = \pm F$, giving
\begin{equation}
\langle F, F \rangle
 = \pm\,\langle F, \star F\rangle
 = 0 .
\end{equation}
By \eqref{eq:trace} the free stress tensor is traceless, for any improvement
coefficient. Consequently $\operatorname{Tr}(\tau)$ contributes nothing at field
degree two and first contributes at degree four.

\paragraph{Scope of the argument.}
The derivation uses only self-duality (or anti-self-duality) and the parity of
the form degree. It does not use the detailed form of the free dynamics, the
choice of improvement term, or the overall normalisation. That robustness is why
we regard the input as settled, and why the remaining question for review is
narrow: whether the intended free theory and stress convention are the ones fixed
in appendix~\ref{app:conventions}.

\paragraph{Computational verification.}
A dedicated test module re-derives the statement from the five-form and the
Hodge star alone, importing none of the flow machinery. It checks:
\begin{enumerate}\itemsep2pt
\item $\star^2 = +\starSquared$, without which the eigenspace decomposition is
      not the one assumed;
\item $\langle F,F\rangle = 0$ on self-dual samples, at eight seeds;
\item $\langle F,F\rangle = 0$ on \emph{anti}-self-dual samples, which the
      argument predicts equally and which would fail if the vanishing were an
      artefact of the projector rather than of the wedge identity;
\item the control: $\langle F,F\rangle \neq 0$ on an unprojected five-form. This
      is the test that gives the other three their content, since a quantity that
      vanished identically would satisfy them all vacuously;
\item the conclusion at four primes;
\item agreement with the production leading-degree table, read once and only to
      be checked against the independent derivation.
\end{enumerate}

\paragraph{A correction worth recording.}
An earlier version of the control called the projector routine with an argument
believed to select the anti-self-dual space, when in fact that argument selects
Lorentzian versus Euclidean signature. The apparent non-vanishing it reported was
a Euclidean projector, where $\star^2 = -1$ and the eigenspace decomposition is
different. Corrected, both eigenspaces vanish exactly as the argument predicts.
The episode is the reason item (1) above is checked explicitly rather than
assumed.

\paragraph{The counterfactual.}
Forcing $\operatorname{Tr}(\tau)$ to begin at degree two activates
$24$ additional targets and raises the degree-ten reachable dimension from
$\dimDtenQ$ to $\dimAten$, giving $\Qten = \gTenCounterfactualQten$ instead of
$\gTenDerivedQten$. This is retained as a regression fixture: it reproduces
precisely the incorrect result this project obtained at an earlier stage, so the
test fails loudly if the assumption is ever silently reverted.
